\documentclass[10pt,a4paper]{amsart}
\usepackage[margin=19mm]{geometry}
\usepackage{amsmath,amssymb,mathtools}
\usepackage[scr=rsfs,cal=euler]{mathalfa}
\usepackage{xfrac}
\usepackage[unicode,hidelinks,bookmarks=false]{hyperref}
\newcommand{\ie}{i.e.\ }
\newcommand{\cf}{cf.\ }
\newcommand{\resp}{respectively\ }
\newcommand{\Subsection}{\subsection}
\providecommand{\nobreakdash}{\nobreak\hbox{-}}
\setlength{\emergencystretch}{2em}
\multlinegap0pt
\usepackage{bm}
\usepackage{bbm}
\usepackage{dsfont}
\usepackage{xspace}
\usepackage{cancel}
\usepackage{mathtools}
\usepackage{amsthm}
\makeatletter
\@ifundefined{theorem}{\newtheorem{theorem}{Theorem}}{}
\@ifundefined{lemma}{\newtheorem{lemma}[theorem]{Lemma}}{}
\@ifundefined{proposition}{\newtheorem{proposition}[theorem]{Proposition}}{}
\makeatother
\title[Singularity Formation: Synergy in Theoretical, Numerical and]{Singularity Formation: Synergy in Theoretical, Numerical and Machine Learning Approaches}
\author{Yixuan Wang}
\date{Source: arXiv:2604.16842v1 (CC BY 4.0)}
\begin{document}
\maketitle

\noindent\emph{Source and licence.} Singularity Formation: Synergy in Theoretical, Numerical and Machine Learning Approaches. Version arXiv:2604.16842v1 (CC BY 4.0), distributed under the Creative Commons Attribution 4.0 International licence, as recorded in the arXiv licence metadata for this exact version and shown on the abstract page https://arxiv.org/abs/2604.16842v1. Full rights notice: licenses/arxiv-2604.16842-CC-BY-4.0.txt.\par\smallskip

\noindent\textit{Editorial scope.} Counted unit: the Proposition whose source label is \texttt{prop d} in the article (source0.tex lines 2347--2349, source label \texttt{prop d}), delivered together with the article's own complete original proof of it (source0.tex lines 2350--2397, one \texttt{proof} environment of 48 lines). No mathematics is written, repaired or re-derived: statement and proof are the article's own text line for line. Required context retained uncounted: dampl (lines 2332--2336); cap (lines 2762--2767). Every definition, display and label of those lines is retained unchanged. Omitted from this package and not redistributed: 1--2331, 2337--2346, 2398--2761, 2768--4117. Nothing inside a retained excerpt is omitted or altered: every retained body is delivered exactly as the article's lines are written, so no omission receipt of any kind accompanies this package. Citations: the retained excerpts carry no citation command at all, so no bibliography entry of the article is transcribed and no reference is redistributed. Reference adaptation: none was needed and none was made; every reference inside the delivered unit resolves inside it, so no unresolved reference or citation is produced. Printed slips kept literal: no printed word, symbol, hypothesis or number of the counted statement or of its proof was altered. Layout and class: the article's own class is \texttt{caltech\_thesis} with packages including url, lipsum, amsthm, mathtools, amssymb, amsmath, graphicx, todonotes, inputenc, fontenc, newtxtext, newtxmath, biblatex; the wrapper is the lane's pinned \texttt{amsart} editorial wrapper at 10pt on A4, which restates the theorem-like environments the retained text needs and adds the article's own macro definitions that the retained text needs (none), with extra wrapper packages (bm, bbm, dsfont, xspace, cancel, mathtools). The delivered numbering is the unit's own. Assets: no retained span contains a figure environment, a graphics-inclusion command, a PDF-page inclusion, a table or a TikZ picture; no image file is copied into this package, the package declares no asset dependency and no image file is redistributed with it. Licence: version arXiv:2604.16842v1 of this article is distributed under the Creative Commons Attribution 4.0 International licence, as recorded in the arXiv licence metadata for this exact version and shown on the abstract page https://arxiv.org/abs/2604.16842v1. A full rights notice is delivered at licenses/arxiv-2604.16842-CC-BY-4.0.txt. Characters: every retained excerpt is pure ASCII, so the delivered engine, XeLaTeX with its default Latin Modern text font, reports no missing character.
\begin{lemma}
\label{dampl}We have the following identity
$$(\sin x f_x,f\rho)=\frac{1}{2}(f^2,\rho)\,,$$
which can be verified directly by using integration by parts.
\end{lemma}

\begin{lemma}
\label{cap}
    Assume $\sum_{k\geq 1} a_k^2+c_k^2<\infty$, then we have the following inequality for \eqref{dample}
    $$\begin{aligned}
 F(a,c)&\coloneqq \sum_{k\geq1}\{ a_k^2(0.84+\frac{1}{k^2}-\frac{1}{(k-1)^2})+c_k^2(0.84+\frac{1}{k(k+1)})+ 2a_k a_{k+1}\frac{1}{(k+1)^2}\\&+2 a_k \sum_{j>k+1}a_j(\frac{1}{j^2}-\frac{1}{(j-1)^2})+2a_kc_k\frac{1+2k-k^2}{2k^2(k+1)}+2a_{k+1}c_k\frac{k^2-k-1}{2k^2(k+1)^2}\\&-2a_{k+2}c_k\frac{k+2}{2(k+1)^2}+\sum_{j>k}2a_k c_j\frac{1}{j(j+1)}\}\geq0\,.\end{aligned}$$
\end{lemma}

\begin{proposition}\label{prop d} The following energy estimate holds for the leading linearized operators
$$dE_1\coloneqq((L_1)_x,u_x\rho)+(L_2,\omega\rho)\leq -0.16[(u_x,u_x\rho)+(\omega,\omega\rho)]\,.$$
\end{proposition}

\begin{proof}
Consider the expansion of $\omega$, $u_x$, and $u$ in the orthonormal basis
$$\omega=\sum_{k\geq 1} a_k o^k\,, \quad u=\sum_{k\geq 1} b_k o^k\,, \quad u_x=\sum_{k\geq1} c_k e^k\,.$$
Note the summation index for $u_x$ satisfies $k \geq 1$ since we can easily see that $$ (u_x,e^0\rho)=\frac{1}{2\pi}\int_0^{2\pi} u_x=0\,.$$We first express $b_k$ in terms of $c_k$. If we insert the expression of the basis into $u$, take derivative and compare the coefficients with the expansion of $u_x$, we get $$c_i=\sum_{k=1}^i b_k-ib_{i+1}\,.$$
Therefore we can solve $$b_{i+1}=b_1-\sum_{k=1}^{i-1}\frac{c_k}{k(k+1)}-\frac{c_{i}}{i}\,.$$
Moreover, we have the compatibility condition $u_x(0)=\sum_{k \geq 0} b_k=0$. Therefore we can solve $b_1$ and obtain 
\begin{equation}\label{bc}
    b_{i}=\sum_{k\geq i}\frac{c_k}{k(k+1)}-\frac{c_{i-1}}{i}\,,
\end{equation}
where we define $c_0=0$.

Now we write out the terms explicitly using the expansions
\begin{equation*}
\begin{aligned}
dE_1&=2(-u\sin x-u_{xx}\sin x,u_x\rho)+2(\sin x\psi+\sin x\psi_{xx},u_x\rho)\\&+2(- \sin x \omega_x- \cos x\psi,\omega \rho) +2 (u \cos{x}+\sin x u_x,\omega \rho)\\&=-[(u_x,u_x\rho)+(\omega,\omega\rho)+(u,u\rho)]+2[-( \cos x\psi,\omega \rho)+(\sin x\psi,u_x\rho)+ (u \cos{x},\omega \rho)]\,.    
\end{aligned}    
\end{equation*}

Here we use the crucial Lemma \ref{dampl} to extract damping on the local terms and the Biot-Savart law $-\psi_{xx}=\omega$ to cancel the effect of the nonlocal terms {$\sin x\psi_{xx}$ in $(L_1)_x$ and $\sin xu_{x}$ in $L_2$}. Next, we calculate the remaining nonlocal terms explicitly.
$$-2( \cos x\psi,\omega \rho)=(2(1-\cos x)\psi,\omega \rho)-2(\psi,\omega \rho)=\frac{1}{\pi}(\psi,\omega)-2(\psi,\omega \rho)\,.$$
We express $\omega$ and $\psi$ both in terms of orthonormal basis $o^k$ corresponding to the weighted norm and the canonical basis $\sin(kx)$ corresponding to the (normalized by $\frac{1}{\pi}$) $L^2$ norm.
$$\omega=\sum_{k \geq 1} (a_k-a_{k+1}) \sin(kx)\,,$$
where we denote $a_0=0$. Therefore 
$$\psi=\sum_{k \geq 1} \frac{a_k-a_{k+1}}{k^2} \sin(kx)\,.$$
Furthermore, we collect
$$\psi=\sum_{k\geq 1} \sum_{j\geq k}\frac{a_j-a_{j+1}}{j^2} o^k\,.$$
Therefore we can compute explicitly that 
$$-2( \cos x\psi,\omega \rho)=\sum_{k \geq 1} \frac{(a_k-a_{k+1})^2}{k^2}-2\sum_{k\geq 1} a_k\sum_{j\geq k}\frac{a_j-a_{j+1}}{j^2}\,.$$

We use integration by parts similar to Lemma \ref{dampl} to obtain 
$$2[(\sin x\psi,u_x\rho)+ (u \cos{x},\omega \rho)]=2(\cos{x}\omega+\psi-\sin x\psi_x,u\rho)\coloneqq2(T_u,u\rho)\,.$$
We further have $$\begin{aligned}
T_u&=\sum_{k \geq 1} (a_k-a_{k+1})[\frac{k+1}{2k}\sin((k-1)x)+\frac{k-1}{2k}\sin((k+1)x)+\frac{1}{k^2}\sin(kx)]\\&=\sum_{k \geq 1} \sin(kx) [\frac{k+2}{2(k+1)}(a_{k+1}-a_{k+2})+\frac{a_{k}-a_{k+1}}{k^2}+\frac{k-2}{2(k-1)}(a_{k-1}-a_{k})]\\&=\sum_{k\geq 1}[\frac{k-2}{2(k-1)}a_{k-1}+(\frac{1}{2k(k-1)}+\frac{1}{k^2})a_{k}+(\frac{k+1}{2k}+\frac{1}{(k+1)^2}-\frac{1}{k^2})a_{k+1}\\&+\sum_{j>k+1}(\frac{1}{j^2}-\frac{1}{(j-1)^2})a_j]o^k\,,  
\end{aligned}$$
where the terms involving $\frac{1}{k-1}$ in the summand is regarded as $0$ for $k=1$. Therefore we collect explicitly that 
$$\begin{aligned}
dE_1&=\sum_{k\geq 1}\{ -(a_k^2+b_k^2+c_k^2)+ \frac{(a_k-a_{k+1})^2}{k^2}-2 a_k\sum_{j\geq k}\frac{a_j-a_{j+1}}{j^2}+2b_k[\frac{k-2}{2(k-1)}a_{k-1}\\&+(\frac{1}{2k(k-1)}+\frac{1}{k^2})a_{k}+(\frac{k+1}{2k}+\frac{1}{(k+1)^2}-\frac{1}{k^2})a_{k+1}+\sum_{j>k+1}(\frac{1}{j^2}-\frac{1}{(j-1)^2})a_j]\}\,.\end{aligned}$$
Substituting \eqref{bc} into the above and we can simplify
$$
\begin{aligned}
dE_1&=-\sum_{k \geq 1}\{ a_k^2(1+\frac{1}{k^2}-\frac{1}{(k-1)^2})+c_k^2(1+\frac{1}{k(k+1)})+ 2a_k a_{k+1}\frac{1}{(k+1)^2}\\&+2 a_k \sum_{j>k+1}a_j(\frac{1}{j^2}-\frac{1}{(j-1)^2})+2a_kc_k\frac{1+2k-k^2}{2k^2(k+1)}+2a_{k+1}c_k\frac{k^2-k-1}{2k^2(k+1)^2}\\&-2a_{k+2}c_k\frac{k+2}{2(k+1)^2}+\sum_{j>k}2a_k c_j\frac{1}{j(j+1)}\}\,.\end{aligned}$$
After this explicit computation, we notice that the damping estimate in Proposition \ref{prop d} can be cast into an estimate of a quadratic form; see \eqref{dample}, which is equivalent to a lower bound on the eigenvalues of an infinite-dimensional symmetric matrix. 
\begin{equation}
    \label{dample}\begin{aligned}
 F(a,c)&\coloneqq \sum_{k\geq1}\{ a_k^2(0.84+\frac{1}{k^2}-\frac{1}{(k-1)^2})+c_k^2(0.84+\frac{1}{k(k+1)})+ 2a_k a_{k+1}\frac{1}{(k+1)^2}\\&+2 a_k \sum_{j>k+1}a_j(\frac{1}{j^2}-\frac{1}{(j-1)^2})+2a_kc_k\frac{1+2k-k^2}{2k^2(k+1)}+2a_{k+1}c_k\frac{k^2-k-1}{2k^2(k+1)^2}\\&-2a_{k+2}c_k\frac{k+2}{2(k+1)^2}+\sum_{j>k}2a_k c_j\frac{1}{j(j+1)}\}\geq0\,.\end{aligned}\end{equation}

We notice that the entries decay fast. {Therefore the strategy to prove \eqref{dample} is to combine a computer-assisted estimate of the eigenvalues of its finite truncation with a decay estimate of the remaining part. We will defer the proof of \eqref{dample} to the Appendix, see Lemma \ref{cap} and the proof}. Thereby we conclude the linear estimate.
\end{proof}

\end{document}
